\documentclass[uplatex,a4paper,11pt]{jsarticle}
\usepackage[margin=22mm]{geometry}
\usepackage[dvipdfmx]{hyperref,xcolor}
\usepackage{fvextra,enumitem,longtable,booktabs,amssymb}
\hypersetup{
  colorlinks=true,
  linkcolor=blue!55!black,
  urlcolor=blue!65!black,
  pdftitle={TeXNote・LilyPondNote P0リバースプロキシ・ファイアウォール構築手順}
}
\definecolor{commandbg}{RGB}{244,246,248}
\DefineVerbatimEnvironment{command}{Verbatim}{
  fontsize=\small,
  frame=single,
  rulecolor=\color{gray!55},
  breaklines=true,
  breakanywhere=true
}
\setlist{nosep}

\title{TeXNote・LilyPondNote P0\
リバースプロキシ・ファイアウォール構築手順\
\large P1認証サーバーへ安全に接続する公開入口}
\author{}
\date{2026年8月31日}

\begin{document}
\maketitle

\begin{abstract}
本書は、TeXNoteとLilyPondNoteの公開入口となるRaspberry Pi「P0」に、Nginx
リバースプロキシとUFWファイアウォールを構築し、P1上の二つの認証サービスへ接続する
手順をまとめたものである。P0には認証処理、ユーザーDB、JWT署名鍵、P2内部トークンを
置かない。P0はHTTPS終端と公開APIの絞り込みだけを担当し、認証とP2への中継はP1が
引き続き担当する。
本文を上から順に実行すると、ポート分離、リバースプロキシ、ファイアウォール、HTTPSを
備えた完成構成になる。
\end{abstract}

\tableofcontents
\newpage

\section{最初に区別すること}

P0とP1は役割が異なる。P0へP1のプログラムをコピーしたり、P0にユーザーDBを作成したり
してはならない。

\begin{longtable}{@{}p{18mm}p{51mm}p{77mm}@{}}
\toprule
機器 & 役割 & 配置するもの\\
\midrule
P0 & インターネット側の公開入口 & Nginx、TLS証明書、UFW、アクセスログ\\
P1 & 認証とAPI中継 & TeX用・LilyPond用FastAPI、JWT署名鍵、SQLite DB、P2内部トークン\\
P2 & 処理サービス & TeX API、LilyPond PDF API、LilyPond移調テキストAPI\\
\bottomrule
\end{longtable}

完成時の通信経路は次のとおりである。

\begin{command}
TeXNote client       -- HTTPS --> P0 -- HTTP --> P1 TeX認証 TCP 8000
                                                   |
                                                   v
                                                P2 TeX TCP 8000

LilyPondNote client  -- HTTPS --> P0 -- HTTP --> P1 Lily認証 TCP 8001
                                                   |-- PDF生成 -> P2 TCP 8001
                                                   `-- 移調テキスト -> P2 TCP 8002
\end{command}

インターネットへ公開するのはP0のTCP 443だけである。P1、P2、TCP 8000、8001、8002、
SSHを直接公開しない。

\section{前提値と置換する値}

本書では完成構成の例として次の値を使う。構築先のネットワークが異なる場合は、
\texttt{P0\_IP}と\texttt{P1\_IP}を実際の固定アドレスへ読み替える。

\begin{longtable}{@{}p{52mm}p{48mm}p{48mm}@{}}
\toprule
項目 & 本書の例 & 実機で確認する場所\\
\midrule
P0の固定IP & \texttt{192.168.3.22} & ルーター、DHCP予約、P0\\
P1（認証）の固定IP & \texttt{192.168.30.10} & ルーター、DHCP予約、P1\\
P2（処理）の固定IP & \texttt{192.168.40.10} & ルーター、DHCP予約、P2\\
管理端末とアプリのLAN & \texttt{192.168.3.0/24} & Mac、iPad、iPhoneのLAN\\
P1 TeX認証API & TCP 8000 & \texttt{texnote-p1.service}\\
P1 LilyPond認証API & TCP 8001 & \texttt{lilypondnote-p1.service}\\
P0 TeXNote公開入口 & TCP 8080 & TeXNoteに入力するURL\\
P0 LilyPondNote公開入口 & TCP 8081 & LilyPondNoteに入力するURL\\
公開ドメイン & \texttt{tex.example.com}、\texttt{lilypond.example.com} & DNS設定\\
HTTP要求上限 & TeX 32 MiB、LilyPond 8 MiB & 各P1実装と一致\\
\bottomrule
\end{longtable}

\texttt{192.168.40.10}は、TeXコンパイル、LilyPond PDF生成、LilyPond移調テキスト生成を
実行する\textbf{P2処理サーバーの固定IP}である。P2は処理用ネットワーク
\texttt{192.168.40.0/24}に置き、インターネットやP0から直接接続させない。P1だけが、
認証・利用上限確認後の内部要求をP2へ送る。

\textbf{重要：}P1とP2は別ホストなので、両ホストで同じポート番号を使用しても衝突しない。
P1ではTeX認証が8000、LilyPond認証が8001、P2ではTeXコンパイルが8000、
LilyPond PDF生成が8001、LilyPond移調テキスト生成が8002である。LilyPondの二つの
サービスはP1上の同じ認証、ユーザーDB、JWTを使用する。

\section{作業前の安全確認}

\begin{enumerate}
  \item P0、P1には固定IPまたはDHCP予約を設定する。
  \item UFWを有効にする前に、管理端末からのSSH許可規則を必ず追加する。
  \item 作業中のSSH接続を残したまま、別ターミナルで新しいSSH接続を試験する。
  \item P1の認証APIとP2への中継が完成していることを確認してからP0を構築する。
  \item ルーターのポート転送は、HTTPSの試験直前まで作成しない。
  \item IPv6を使う場合はIPv6にも同じ境界制御を適用する。できない場合はサーバー用
        ネットワークへのグローバルIPv6配布を停止する。
\end{enumerate}

P1の稼働状態と実際の待受を記録する。

\begin{command}
# P1で実行
systemctl status lilypondnote-p1.service --no-pager -l
systemctl is-enabled lilypondnote-p1.service
sudo ss -ltnp | grep -E ':8000|:8001'
\end{command}

\section{P0のOSとネットワークを準備する}

64-bit Raspberry Pi OS LiteまたはDebianを導入する。管理ユーザーの公開鍵SSHを設定し、
パスワードだけのSSHログインは避ける。P0の固定IPを設定した後、P0で更新する。

\begin{command}
sudo apt update
sudo apt full-upgrade
sudo reboot
\end{command}

再接続後、アドレスとP1への経路を確認する。

\begin{command}
ip -br address
ip route
ping -c 3 192.168.30.10
\end{command}

この時点ではpingに応答しなくてもよいが、最終的にP0からP1のTCP 8000と8001へ到達できる
ルーティングとACLが必要である。

\section{P0にNginxを構築する}

\subsection{導入と初期設定の保全}

\begin{command}
sudo apt install nginx
nginx -v
sudo cp -a /etc/nginx /etc/nginx.before-lilypondnote
systemctl status nginx --no-pager -l
\end{command}

\subsection{P0用サイト設定を作る}

P0で次のファイルを\texttt{/etc/nginx/sites-available/compile-services-p0}として作成する。
ローカルDNSを設けないLAN内試験では、同じP0のIPアドレスを使い、TeXNoteはTCP 8080、
LilyPondNoteはTCP 8081で入口を分ける。したがって\texttt{server\_name}による振分けには
依存しない。
TeX側は\texttt{/login}、\texttt{/me}、\texttt{/compile}だけを転送する。LilyPond側は
これらに\texttt{/transpose}を加えた4パスだけを転送し、FastAPIの\texttt{/docs}や
\texttt{/openapi.json}を公開しない。

\begin{command}
server {
    listen 8080 default_server;
    listen [::]:8080 default_server;
    server_name _;
    server_tokens off;
    client_max_body_size 32m;
    client_body_timeout 30s;
    access_log /var/log/nginx/texnote-p0-access.log;
    error_log  /var/log/nginx/texnote-p0-error.log;
    location = /login {
        proxy_pass http://192.168.30.10:8000;
        include /etc/nginx/proxy_params;
    }
    location = /me {
        proxy_pass http://192.168.30.10:8000;
        include /etc/nginx/proxy_params;
    }
    location = /compile {
        proxy_pass http://192.168.30.10:8000;
        include /etc/nginx/proxy_params;
        proxy_connect_timeout 5s;
        proxy_send_timeout 120s;
        proxy_read_timeout 120s;
    }
    location / {
        return 404;
    }
}

server {
    listen 8081 default_server;
    listen [::]:8081 default_server;
    server_name _;
    server_tokens off;
    client_max_body_size 8m;
    client_body_timeout 30s;
    access_log /var/log/nginx/lilypondnote-p0-access.log;
    error_log  /var/log/nginx/lilypondnote-p0-error.log;
    location = /login {
        proxy_pass http://192.168.30.10:8001;
        include /etc/nginx/proxy_params;
    }
    location = /me {
        proxy_pass http://192.168.30.10:8001;
        include /etc/nginx/proxy_params;
    }
    location = /compile {
        proxy_pass http://192.168.30.10:8001;
        include /etc/nginx/proxy_params;
        proxy_connect_timeout 5s;
        proxy_send_timeout 60s;
        proxy_read_timeout 60s;
    }
    location = /transpose {
        proxy_pass http://192.168.30.10:8001;
        include /etc/nginx/proxy_params;
        proxy_connect_timeout 5s;
        proxy_send_timeout 30s;
        proxy_read_timeout 30s;
    }
    location / {
        return 404;
    }
}

server {
    listen 80 default_server;
    listen [::]:80 default_server;
    server_name _;
    return 444;
}
\end{command}

TCP 80の最後のブロックは、ポート番号を省略した要求を認証APIへ流さず切断する。
TeXNote側とLilyPondNote側の両方へ\texttt{listen 80}や
\texttt{server\_name 192.168.3.22}を書いてはならない。同じIPアドレスとTCP 80の
\texttt{server\_name}が競合し、一方の設定が無視されるためである。LAN内のクライアントは、
TeXNoteに\texttt{http://192.168.3.22:8080}、LilyPondNoteに
\texttt{http://192.168.3.22:8081}を設定する。

\texttt{proxy\_pass}に末尾のスラッシュを付けず、クライアントから受け取った
TeX側の3パスとLilyPond側の4パスをそのままP1へ渡す。P0は\texttt{/transpose}の
本文を解釈・変換せず、認証も行わない。

設定を有効化する。

\begin{command}
sudo ln -s /etc/nginx/sites-available/compile-services-p0 \
  /etc/nginx/sites-enabled/compile-services-p0
sudo unlink /etc/nginx/sites-enabled/default
sudo nginx -t
sudo systemctl reload nginx
sudo systemctl enable nginx
sudo ss -ltnp | grep -E ':80|:8080|:8081'
\end{command}

\texttt{nginx -t}に\texttt{conflicting server name}が表示されず、\texttt{ss}に
TCP 8080と8081のNginx待受が表示されることを確認する。

\section{P0にUFWを構築する}

\subsection{許可規則を先に登録する}

P0で実行する。管理LAN以外からSSHを許可しない。Mac、iPad、iPhoneが属する
\texttt{192.168.3.0/24}からTCP 8080と8081を許可する。TCP 80は構築試験中だけ
同じLANから許可する。TCP 443は公開サービスなので全送信元から許可する。

\begin{command}
sudo apt install ufw
sudo ufw default deny incoming
sudo ufw default allow outgoing
sudo ufw allow from 192.168.3.0/24 to any port 22 proto tcp \
  comment 'SSH from management network'
sudo ufw allow from 192.168.3.0/24 to any port 80 proto tcp \
  comment 'temporary HTTP test'
sudo ufw allow from 192.168.3.0/24 to any port 8080 proto tcp \
  comment 'TeXNote P0 LAN entry'
sudo ufw allow from 192.168.3.0/24 to any port 8081 proto tcp \
  comment 'LilyPondNote P0 LAN entry'
sudo ufw allow 443/tcp comment 'LilyPondNote HTTPS'
sudo ufw show added
\end{command}

表示されたSSH許可規則が正しいことを確認してから有効化する。

\begin{command}
sudo ufw enable
sudo ufw status verbose
sudo ufw status numbered
\end{command}

番号付き一覧に、\texttt{192.168.3.0/24}からTCP 8080と8081への\texttt{ALLOW IN}が
各1行表示されることを確認する。Nginxが待ち受けていても、この二規則がなければMacからの
接続はタイムアウトする。

別ターミナルからP0へ新しくSSH接続する。成功するまで元のSSH接続を閉じない。

\subsection{P0の送信制限について}

上記の\texttt{default allow outgoing}では、P0自身のUFWは送信先を制限しない。
そのため、VLANルーターまたはL3スイッチのACLで次を強制する。

\begin{longtable}{@{}p{34mm}p{40mm}p{30mm}p{38mm}@{}}
\toprule
送信元 & 宛先 & ポート & 動作\\
\midrule
P0 & P1 & TCP 8000、8001 & 許可\\
P0 & P2 & すべて & 拒否\\
DMZ（P0） & 認証VLAN & TCP 8000、8001以外 & 拒否\\
管理ネットワーク & P0 & TCP 22 & 許可\\
Internet & P0 & TCP 443 & 許可\\
Internet & P0/P1/P2 & 上記以外 & 拒否\\
\bottomrule
\end{longtable}

OS更新に必要なP0からのDNS、NTP、HTTPSは、利用するルーターの仕様に合わせて許可する。
ホスト側でも厳格な送信制限を行う場合は、DNSサーバー、NTPサーバー、更新先の扱いを
決めてからUFWのoutgoing規則を変更する。未確認のまま\texttt{deny outgoing}にすると、
名前解決、時刻同期、証明書更新を止める危険がある。

\section{P1をP0からだけ受け付けるようにする}

この節だけはP0ではなく\textbf{P1で実行する}。詳細はP1構築手順書に従う。
TeX認証サービスはTCP 8000、LilyPond認証サービスはTCP 8001を使用する。

\subsection{P1の待受を確認する}

P1とP0が別サブネットの場合、P1は\texttt{127.0.0.1}ではなくP1の認証ネットワーク側
固定IPで待ち受ける。二つのsystemdサービスを次のように設定する。

\begin{command}
# /etc/systemd/system/lilypondnote-p1.service のExecStart例
# texnote-p1.service
ExecStart=/opt/texnote/P1/.venv/bin/uvicorn main:app \
  --host 192.168.30.10 --port 8000

# lilypondnote-p1.service
ExecStart=/opt/lilypondnote/P1/.venv/bin/uvicorn main:app \
  --host 192.168.30.10 --port 8001
\end{command}

編集後に反映する。

\begin{command}
sudo systemctl daemon-reload
sudo systemctl restart texnote-p1.service lilypondnote-p1.service
systemctl status texnote-p1.service --no-pager -l
systemctl status lilypondnote-p1.service --no-pager -l
sudo ss -ltnp | grep -E ':8000|:8001'
\end{command}

\subsection{P1のUFWでP0だけを許可する}

SSH許可を先に登録し、P0の固定IPからTCP 8000と8001への通信だけを許可する。

\begin{command}
sudo apt install ufw
sudo ufw default deny incoming
sudo ufw default allow outgoing
sudo ufw allow from 192.168.3.0/24 to any port 22 proto tcp \
  comment 'SSH from management network'
sudo ufw allow from 192.168.3.22 to any port 8000 proto tcp \
  comment 'P0 to TeX authentication'
sudo ufw allow from 192.168.3.22 to any port 8001 proto tcp \
  comment 'P0 to LilyPond authentication'
sudo ufw show added
sudo ufw enable
sudo ufw status verbose
\end{command}

既にUFWが有効な場合は、\texttt{sudo ufw status numbered}で規則を確認する。
TCP 8000または8001をLAN全体や全送信元へ許可する古い規則があれば、P0からの接続成功を
確認した後、表示された番号を指定して削除する。

\begin{command}
sudo ufw status numbered
sudo ufw delete <古い全面許可規則の番号>
\end{command}

P1からP2へのURLはサービス別に設定する。ここで使用する\texttt{192.168.40.10}は、
「前提値と置換する値」で示したP2処理サーバーである。P2上のサービスは次の3個すべてを
設定する。

\begin{longtable}{@{}p{43mm}p{28mm}p{68mm}@{}}
\toprule
処理 & P2待受 & P1の環境変数\\
\midrule
TeX PDF生成 & TCP 8000 & \texttt{TEXNOTE\_P2\_TYPESET\_URL}\\
LilyPond PDF生成 & TCP 8001 & \texttt{LILYPONDNOTE\_P2\_TYPESET\_URL}\\
LilyPond移調テキスト生成 & TCP 8002 & \texttt{LILYPONDNOTE\_P2\_TRANSPOSE\_URL}\\
\bottomrule
\end{longtable}

\begin{command}
# /etc/texnote/p1.env
TEXNOTE_P2_TYPESET_URL=http://192.168.40.10:8000/v1/typeset

# /etc/lilypondnote/p1.env
LILYPONDNOTE_P2_TYPESET_URL=http://192.168.40.10:8001/v1/typeset
LILYPONDNOTE_P2_TRANSPOSE_URL=http://192.168.40.10:8002/v1/transpose
\end{command}

PDF生成APIの現行パスは\texttt{/v1/typeset}、移調テキストAPIは
\texttt{/v1/transpose}である。TeX用P2内部トークンとLilyPond用P2内部トークンは別の値に
する。LilyPondのPDFと移調は同じLilyPond認証を通るため、LilyPond用内部トークンを共有する。
環境ファイルを変更した場合は該当するP1サービスを再起動する。

\section{LAN内でP0経由を試験する}

最初に、P0が各ポートの公開対象外パスを404にすることを管理端末から確認する。
本書のP0固定IP \texttt{192.168.3.22}を使って確認する。

\begin{command}
curl -i http://192.168.3.22:8080/openapi.json
curl -i http://192.168.3.22:8081/openapi.json
\end{command}

いずれもHTTP 404でなければ公開を進めない。次に両サービスへログインする。

\begin{command}
curl -sS -X POST http://192.168.3.22:8080/login \
  -H 'Content-Type: application/json' \
  -d '{"email":"tex-user@example.com","password":"REPLACE_ME","service":"texnote"}'

curl -sS -X POST http://192.168.3.22:8081/login \
  -H 'Content-Type: application/json' \
  -d '{"email":"music-user@example.com","password":"REPLACE_ME","service":"lilypondnote"}'
\end{command}

各JWTは発行元サービスだけで使用する。TeX JWTをLilyPondへ、LilyPond JWTをTeXへ
送った場合に認証が失敗することを確認する。トークンをシェル履歴へ残さない。

\begin{command}
TEX_TOKEN='TeXサービスから取得したaccess_token'
curl -sS http://192.168.3.22:8080/me \
  -H "Authorization: Bearer $TEX_TOKEN"
\end{command}

\texttt{/me}のJSONに\texttt{"service":"texnote"}が含まれることを確認する。
LilyPondNote側はTCP 8081の\texttt{/me}が\texttt{"service":"lilypondnote"}を返す。
また、TeXNoteのJSONをTCP 8081へ送る試験、およびLilyPondNoteのJSONをTCP 8080へ送る
試験はHTTP 422となり、サービスを取り違えたログインが成立しないことを確認する。

最後にTeXNoteのURLを\texttt{http://P0のIP:8080}、LilyPondNoteのURLを
\texttt{http://P0のIP:8081}へ設定し、
ログインと実際のコンパイルを確認する。
問題がある場合は次のログを同時に調べる。

\begin{command}
# P0
sudo tail -f /var/log/nginx/texnote-p0-access.log \
  /var/log/nginx/lilypondnote-p0-access.log \
  /var/log/nginx/texnote-p0-error.log \
  /var/log/nginx/lilypondnote-p0-error.log

# P1
journalctl -u texnote-p1.service -f
journalctl -u lilypondnote-p1.service -f
\end{command}

さらに、ネットワーク境界が設計どおりかを次の4方向から確認する。TCP接続試験には
\texttt{nc}を使う。Raspberry Pi OSでコマンドがない場合は、試験するP0/P1へ
\texttt{sudo apt install netcat-openbsd}を導入する。成功時は通常\texttt{succeeded}、
拒否時は\texttt{refused}または\texttt{timed out}と表示される。

\subsection{P0からP1の8000・8001へ接続できること}

\textbf{P0で実行する。}P0から二つのP1認証APIへ到達できることを確認する。

\begin{command}
nc -vz -w 3 192.168.30.10 8000
nc -vz -w 3 192.168.30.10 8001
\end{command}

両方が\texttt{succeeded}なら正常である。失敗する場合は、P1のsystemd待受IP、P1のUFW、
P0--P1間のルーターACLを確認する。HTTPレベルも確認する場合は、P0から
\texttt{curl -i http://192.168.30.10:8000/docs}等を実行し、200、401、404、405のいずれかが
返ればTCP/HTTPは到達している。タイムアウトや接続拒否は未到達である。

\subsection{P0からP2の8000・8001・8002へ直接接続できないこと}

\textbf{P0で実行する。}P0はP1へだけ接続し、P2へ直接到達してはならない。

\begin{command}
nc -vz -w 3 192.168.40.10 8000
nc -vz -w 3 192.168.40.10 8001
nc -vz -w 3 192.168.40.10 8002
\end{command}

3個すべてが\texttt{refused}または\texttt{timed out}なら境界は正常である。一つでも
\texttt{succeeded}なら、DMZから処理VLANへのルーターACLが広すぎるため、P0からP2への
通信を拒否する。P2サービスを一時停止して「失敗」に見せるのではなく、P2が稼働中の状態で
ACLによって遮断されることを確認する。

\subsection{管理端末・一般LANからP1へ直接接続できないこと}

\textbf{管理用Macまたは一般LAN端末で実行する。}P0を経由しないP1認証APIへの接続が
遮断されることを確認する。SSH用管理VLANからの試験でも、8000・8001は許可しない。

\begin{command}
nc -vz -w 3 192.168.30.10 8000
nc -vz -w 3 192.168.30.10 8001
\end{command}

両方が\texttt{refused}または\texttt{timed out}なら正常である。一つでも
\texttt{succeeded}なら、P1のUFWにLAN全体または全送信元からの古い許可規則がないか
\texttt{sudo ufw status numbered}で確認する。P1の8000・8001を許可する送信元はP0の固定IP
\texttt{192.168.3.22}だけにする。

\subsection{P1からP2の8000・8001・8002へ接続できること}

\textbf{P1で実行する。}認証後の要求を全処理サービスへ送れることを確認する。

\begin{command}
nc -vz -w 3 192.168.40.10 8000
nc -vz -w 3 192.168.40.10 8001
nc -vz -w 3 192.168.40.10 8002
\end{command}

3個すべてが\texttt{succeeded}ならネットワーク到達性は正常である。失敗したポートがあれば、
対応するP2サービスの待受、P2のUFW、認証VLANから処理VLANへのACLを確認する。これはTCP
到達性だけの試験なので、続けて実際のTeX PDF生成、LilyPond PDF生成、移調テキスト生成を
P0経由で行い、内部トークンとAPIパスも正しいことを確認する。

\section{HTTPS証明書と外部公開}

\subsection{公開前提}

所有するドメインのDNSをP0へ到達するグローバルIPへ向ける。契約回線がCGNATの場合、
通常のIPv4ポート転送は利用できないことがある。ルーターでインターネットから転送する
のは次だけとする。

\begin{command}
Internet TCP 443 -> P0 192.168.3.22 TCP 443
\end{command}

TCP 22、80、8000、8001を恒常的に転送しない。HTTP-01方式で証明書を初回取得する場合は
TCP 80が一時的に必要になることがある。その場合だけP0へ転送し、取得後に削除する。

\subsection{Certbotで証明書を設定する}

P0で実行する。ドメインは実際の値へ置き換える。

\begin{command}
sudo apt update
sudo apt install certbot python3-certbot-nginx
sudo certbot --nginx -d tex.example.com -d lilypond.example.com
sudo nginx -t
\end{command}

Certbotが作成したNginx設定を確認し、HTTPからHTTPSへのリダイレクトとTCP 443の
\texttt{ssl}待受が存在することを確認する。自動更新を試験する。

\begin{command}
systemctl status certbot.timer --no-pager
sudo certbot renew --dry-run
\end{command}

証明書取得に一時許可したTCP 80を閉じる。ルーター側の一時転送も削除する。

\begin{command}
sudo ufw status numbered
sudo ufw delete <temporary HTTP test規則の番号>
sudo ufw status verbose
\end{command}

外部回線（携帯電話回線など、家庭内Wi-Fi以外）から確認する。

\begin{command}
curl -i https://tex.example.com/openapi.json
curl -i https://lilypond.example.com/openapi.json
\end{command}

両方が404であることを確認する。TeXNoteは\path{https://tex.example.com}、
LilyPondNoteは\path{https://lilypond.example.com}を最終URLにする。TLS証明書エラーを
無視する設定はアプリにもP0にも追加しない。

\section{再起動試験と完了チェックリスト}

P0、P1、P2を一台ずつ再起動し、各回でログインとコンパイルが復旧することを確認する。

\begin{command}
# 各サーバーで該当サービスを確認
systemctl is-enabled nginx
systemctl is-active nginx
sudo ufw status verbose
\end{command}

\begin{itemize}
  \item[$\square$] P0にはNginx、TLS証明書、UFWだけがあり、両サービスの秘密値やDBがない。
  \item[$\square$] P0はHTTPS TCP 443だけを外部公開している。
  \item[$\square$] 二つのドメインをP1の別ポートへ正しく振り分ける。
  \item[$\square$] LilyPond側だけが\texttt{/transpose}をP1へ転送する。
  \item[$\square$] \texttt{/docs}、\texttt{/openapi.json}、その他のパスは404になる。
  \item[$\square$] P1のTCP 8000と8001はP0の固定IPからだけ許可される。
  \item[$\square$] P2のTCP 8000、8001、8002はP1の固定IPからだけ許可される。
  \item[$\square$] P0からP2へ直接接続できない。
  \item[$\square$] SSHは管理ネットワークからだけ許可される。
  \item[$\square$] IPv6にも同等の制限があるか、サーバー側のグローバルIPv6が無効である。
  \item[$\square$] Certbotの更新試験が成功する。
  \item[$\square$] 各サーバーの再起動後にログインとコンパイルが復旧する。
\end{itemize}

\section{緊急停止}

P0に問題が起きた場合は、まずルーターのTCP 443転送を停止して外部公開を止める。
続いてNginxを停止する。P0はDBを持たないため、この操作でP1とP2のデータは変更されない。

\begin{command}
# P0
sudo systemctl stop nginx
\end{command}

原因を確認してP0の設定を修正した後、\texttt{nginx -t}が成功してからNginxと
ルーターのTCP 443転送を順に再開する。障害対応のためにP1を一般LANやインターネットへ
直接公開してはならない。

\section{まとめ}

P0はP1の代替サーバーではなく、P1の前に置く小さな公開境界である。認証情報をP0へ
移さず、LAN内ではTCP 8080と8081、外部公開では二つのドメインでサービスを振り分け、
LilyPond側だけが移調パスを公開する。
P1のUFWではP0の固定IPだけを許可する。
さらにルーターまたはL3スイッチでP0からP2への直接通信を拒否することで、公開入口、
認証、コンパイルの役割を明確に分離できる。

\end{document}
